\documentclass[11pt,a4paper]{article}
\usepackage[utf8]{inputenc}
\usepackage[T1]{fontenc}
\usepackage{lmodern}
\usepackage[a4paper,margin=1in]{geometry}
\usepackage{longtable}
\usepackage{array}
\usepackage{booktabs}
\usepackage{enumitem}
\usepackage{xcolor}
\usepackage{hyperref}
\usepackage{listings}

\hypersetup{
  colorlinks=true,
  linkcolor=blue,
  urlcolor=blue
}

\lstdefinestyle{doccode}{
  basicstyle=\ttfamily\small,
  breaklines=true,
  columns=fullflexible,
  keepspaces=true,
  frame=single,
  rulecolor=\color{black!20}
}

\newcommand{\docpath}[1]{\texttt{#1}}
\newcolumntype{L}[1]{>{\raggedright\arraybackslash}p{#1}}

\title{Anhoc Technical Design}
\date{}

\begin{document}
\maketitle
\tableofcontents
\newpage

\section{Purpose}

This document describes the current technical design of the Anhoc project based on the implemented code in \docpath{frontend/} and \docpath{backend/}.

It covers:
\begin{itemize}[leftmargin=1.5em]
  \item Architecture design
  \item Functional design
  \item Data model overview
  \item REST API specification
  \item Main sequence flows
\end{itemize}

\section{System Overview}

Anhoc is a math learning platform with:
\begin{itemize}[leftmargin=1.5em]
  \item Student learning content by subject and grade
  \item Practice sessions generated from question templates
  \item Grade-level test sessions
  \item Progress tracking through mastery, XP, levels, and achievements
  \item Admin functions for lessons, question templates, users, roles, and access control
\end{itemize}

\section{Architecture Design}

\subsection{Logical Architecture}

The runtime flow is:
\begin{enumerate}[leftmargin=1.5em]
  \item User Browser sends requests to the Next.js frontend.
  \item The frontend calls the Express backend through REST endpoints under \docpath{/api/v1}.
  \item The backend uses Prisma ORM to access PostgreSQL.
  \item The backend also invokes the Math Generation Service, Mastery Service, Level Service, and Achievement Service.
\end{enumerate}

\subsection{Deployment View}

The deployment arrangement is:
\begin{enumerate}[leftmargin=1.5em]
  \item Vercel hosts the frontend runtime.
  \item Render hosts the Node.js API.
  \item Neon PostgreSQL provides the database.
\end{enumerate}

\subsection{Backend Structure}

\begin{itemize}[leftmargin=1.5em]
  \item \docpath{backend/server.ts}
  \begin{itemize}[leftmargin=2em]
    \item Bootstraps Express
    \item Configures CORS and JSON middleware
    \item Mounts versioned APIs under \docpath{/api/v1}
    \item Seeds achievements on startup
  \end{itemize}
  \item \docpath{backend/routes/}
  \begin{itemize}[leftmargin=2em]
    \item \docpath{auth.ts}: login, register, profile, password, activity
    \item \docpath{lessons.ts}: lessons, grades, subjects, practice start, study tracking
    \item \docpath{tests.ts}: attempts, answer submission, template management, grade tests, question reports
    \item \docpath{admin.ts}: users, roles, access control, dashboard statistics
    \item \docpath{achievements.ts}: achievement catalog and award checks
  \end{itemize}
  \item \docpath{backend/middleware/auth.ts}
  \begin{itemize}[leftmargin=2em]
    \item JWT authentication
    \item Permission-based authorization
    \item Self-or-admin access checks
  \end{itemize}
  \item \docpath{backend/services/}
  \begin{itemize}[leftmargin=2em]
    \item \docpath{mathService.ts}: variable generation, answer evaluation, template rendering
    \item \docpath{masteryService.ts}: lesson mastery and study-time updates
    \item \docpath{levelService.ts}: XP accumulation and level calculation
    \item \docpath{achievementService.ts}: seed and award achievements
  \end{itemize}
\end{itemize}

\subsection{Frontend Structure}

\begin{itemize}[leftmargin=1.5em]
  \item \docpath{frontend/src/app/}: App Router pages for auth, student, and admin areas
  \item \docpath{frontend/src/services/}: Axios-based API clients
  \item \docpath{frontend/src/redux/}: Auth and permission state
  \item \docpath{frontend/src/components/}: Student and admin UI components
\end{itemize}

\subsection{Security Model}

\begin{itemize}[leftmargin=1.5em]
  \item Authentication uses JWT in \docpath{Authorization: Bearer <token>}
  \item Frontend stores the token in \docpath{localStorage} and a browser cookie for route middleware or proxy handling
  \item Authorization is based on role name in JWT, grouped permissions per resource, and subject-level restrictions via \docpath{role\_subject\_permissions}
\end{itemize}

\subsection{Core Architectural Decisions}

\begin{itemize}[leftmargin=1.5em]
  \item REST API rather than server actions or GraphQL
  \item Prisma as the data access layer
  \item Question instances are materialized as \docpath{question\_snapshots}
  \item Practice and test completion triggers secondary progression services:
  \begin{itemize}[leftmargin=2em]
    \item mastery update
    \item XP grant
    \item student stats update
    \item achievement evaluation
  \end{itemize}
\end{itemize}

\section{Functional Design}

\subsection{Authentication and Identity}

Capabilities:
\begin{itemize}[leftmargin=1.5em]
  \item Register a new user
  \item Login and receive JWT
  \item Read current profile
  \item Change password
  \item View 7-day XP activity
\end{itemize}

Business behavior:
\begin{itemize}[leftmargin=1.5em]
  \item Registration creates both \docpath{users} and \docpath{student\_stats}
  \item Login expands role permissions into a grouped structure such as:
\end{itemize}

\begin{lstlisting}[style=doccode]
lesson -> [read, manage]
user   -> [manage]
\end{lstlisting}

\subsection{Learning Content Management}

Capabilities:
\begin{itemize}[leftmargin=1.5em]
  \item Manage subjects
  \item Manage grades
  \item Create, update, delete lessons
  \item List lessons with grade and subject metadata
\end{itemize}

Business behavior:
\begin{itemize}[leftmargin=1.5em]
  \item Lessons belong to one grade and one subject
  \item Non-admin access can be restricted to specific subjects
  \item Lesson content is bilingual with \docpath{title\_en}, \docpath{title\_vi}, \docpath{content\_markdown\_en}, and \docpath{content\_markdown\_vi}
\end{itemize}

\subsection{Practice Engine}

Capabilities:
\begin{itemize}[leftmargin=1.5em]
  \item List lessons with available practice
  \item Start lesson-based practice
  \item Answer generated questions
  \item Finish attempt and compute score
  \item View attempt detail and practice history
\end{itemize}

Business behavior:
\begin{itemize}[leftmargin=1.5em]
  \item Practice sessions are stored as \docpath{test\_attempts} with \docpath{is\_practice = true}
  \item Each session creates multiple \docpath{question\_snapshots}
  \item Snapshot data stores template reference, generated variables, accepted right answers, student answer, correctness, and points
  \item At finish:
  \begin{itemize}[leftmargin=2em]
    \item unanswered snapshots are marked incorrect
    \item score is calculated as \docpath{(correct / total) * 100}
    \item lesson mastery is updated for lesson-based practice
    \item XP is awarded as \docpath{correct\_answers * 2}
    \item achievements are re-evaluated
  \end{itemize}
\end{itemize}

\subsection{Grade Test Engine}

Capabilities:
\begin{itemize}[leftmargin=1.5em]
  \item Discover grade tests derived from available lesson templates
  \item Start a 50-question grade test
  \item Submit answers
  \item Finish and score the test
\end{itemize}

Business behavior:
\begin{itemize}[leftmargin=1.5em]
  \item Grade tests are assembled dynamically from all templates under lessons in the selected grade
  \item Grade test attempts use \docpath{is\_practice = false} and \docpath{lesson\_id = null}
  \item The API returns attempts with fully expanded snapshot and lesson metadata
\end{itemize}

\subsection{Question Template Management}

Capabilities:
\begin{itemize}[leftmargin=1.5em]
  \item Create one or many templates
  \item Read, update, delete templates
  \item Report problematic generated questions
  \item Review and resolve reports
\end{itemize}

Business behavior:
\begin{itemize}[leftmargin=1.5em]
  \item Templates support bilingual prompt and explanation content
  \item Difficulty is normalized to \docpath{easy | medium | hard}
  \item Theoretical questions must have exactly 1 correct answer and exactly 3 false answers in \docpath{logic\_config.false\_answers}
  \item Dynamic math questions support variable ranges, choices, derived expressions, constraints, and multiple accepted formulas
\end{itemize}

\subsection{Progression and Gamification}

Capabilities:
\begin{itemize}[leftmargin=1.5em]
  \item Track lesson study time
  \item Track lesson mastery
  \item Award XP and levels
  \item Award achievements
  \item Show achievement catalog and earned state
\end{itemize}

Business behavior:
\begin{itemize}[leftmargin=1.5em]
  \item Mastery status becomes \docpath{not\_started}, \docpath{in\_progress}, or \docpath{completed} when score is at least 80
  \item XP is stored in \docpath{xp\_logs}
  \item Level is recalculated from cumulative XP
  \item Achievements are seeded on startup and checked after completion flows
\end{itemize}

\subsection{Administration}

Capabilities:
\begin{itemize}[leftmargin=1.5em]
  \item List users and roles
  \item Create, update, delete users
  \item Assign role to a user
  \item Manage role permissions and subject permissions
  \item View user performance insights
  \item View dashboard statistics
\end{itemize}

Business behavior:
\begin{itemize}[leftmargin=1.5em]
  \item Access control is split into action or resource permissions and subject-level access
  \item Admin dashboard aggregates users, lessons, attempts, average score, top users, recent activity, and 14-day history
\end{itemize}

\section{Data Model Overview}

\subsection{Main Entities}

\begin{longtable}{L{0.28\textwidth}L{0.62\textwidth}}
\toprule
\textbf{Entity} & \textbf{Purpose} \\
\midrule
\endhead
\docpath{users} & Application users \\
\docpath{roles} & Named roles such as admin or student \\
\docpath{actions}, \docpath{resources}, \docpath{role\_permissions} & Dynamic RBAC \\
\docpath{subjects}, \docpath{grades}, \docpath{lessons} & Learning content hierarchy \\
\docpath{question\_templates} & Reusable question definitions \\
\docpath{test\_attempts} & Practice or test session header \\
\docpath{question\_snapshots} & Materialized generated questions per attempt \\
\docpath{user\_lesson\_mastery} & User progress per lesson \\
\docpath{student\_stats} & Aggregated XP, level, score, activity \\
\docpath{achievements}, \docpath{user\_achievements} & Gamification rewards \\
\docpath{xp\_logs} & XP history \\
\docpath{question\_template\_reports} & Problem report workflow \\
\docpath{audit\_logs}, \docpath{activity\_evidence} & Traceability and user evidence \\
\bottomrule
\end{longtable}

\subsection{Content Hierarchy}

The main relationships are:
\begin{itemize}[leftmargin=1.5em]
  \item \docpath{SUBJECT} contains \docpath{GRADE}
  \item \docpath{SUBJECT} groups \docpath{LESSON}
  \item \docpath{GRADE} contains \docpath{LESSON}
  \item \docpath{LESSON} has \docpath{QUESTION\_TEMPLATE}
  \item \docpath{USER} performs \docpath{TEST\_ATTEMPT}
  \item \docpath{TEST\_ATTEMPT} contains \docpath{QUESTION\_SNAPSHOT}
  \item \docpath{USER} tracks \docpath{USER\_LESSON\_MASTERY}
  \item \docpath{LESSON} is referenced by \docpath{USER\_LESSON\_MASTERY}
\end{itemize}

\subsection{Important Persistence Rules}

\begin{itemize}[leftmargin=1.5em]
  \item A generated practice or test question is not computed on every render; it is persisted as a snapshot
  \item Lesson-based practice attempts retain \docpath{lesson\_id}
  \item Grade tests intentionally do not bind the attempt to one lesson
  \item Student progress is a combination of raw attempts, derived mastery, cumulative XP or level, and achievements
\end{itemize}

\section{API Specification}

Base URL:
\begin{lstlisting}[style=doccode]
/api/v1
\end{lstlisting}

Authentication:
\begin{itemize}[leftmargin=1.5em]
  \item Protected endpoints require \docpath{Authorization: Bearer <token>}
\end{itemize}

Response conventions:
\begin{itemize}[leftmargin=1.5em]
  \item Success returns JSON payloads
  \item Validation or authorization errors generally return \docpath{400}, \docpath{401}, \docpath{403}, \docpath{404}, or \docpath{409}
  \item Unhandled failures return \docpath{500}
\end{itemize}

\subsection{Health}

\begin{longtable}{L{0.12\textwidth}L{0.25\textwidth}L{0.12\textwidth}L{0.41\textwidth}}
\toprule
\textbf{Method} & \textbf{Path} & \textbf{Auth} & \textbf{Description} \\
\midrule
\endhead
\docpath{GET} & \docpath{/health} & No & API health check \\
\bottomrule
\end{longtable}

\subsection{Auth API}

\begin{longtable}{L{0.12\textwidth}L{0.31\textwidth}L{0.12\textwidth}L{0.35\textwidth}}
\toprule
\textbf{Method} & \textbf{Path} & \textbf{Auth} & \textbf{Description} \\
\midrule
\endhead
\docpath{POST} & \docpath{/auth/register} & No & Register user \\
\docpath{POST} & \docpath{/auth/login} & No & Login and return token \\
\docpath{GET} & \docpath{/auth/permissions/:userId} & No & Get grouped permissions by user \\
\docpath{GET} & \docpath{/auth/profile} & Yes & Get current user profile \\
\docpath{PATCH} & \docpath{/auth/password} & Yes & Change current user password \\
\docpath{GET} & \docpath{/auth/activity} & Yes & Get 7-day XP activity \\
\bottomrule
\end{longtable}

Example login request:
\begin{lstlisting}[style=doccode]
{
  "email": "student@example.com",
  "password": "secret123"
}
\end{lstlisting}

Example login response:
\begin{lstlisting}[style=doccode]
{
  "token": "jwt-token",
  "user": {
    "id": "uuid",
    "username": "student1",
    "role": "student",
    "permissions": {
      "lesson": ["read"],
      "test": ["read"]
    }
  }
}
\end{lstlisting}

\subsection{Lessons API}

\begin{longtable}{L{0.12\textwidth}L{0.31\textwidth}L{0.12\textwidth}L{0.35\textwidth}}
\toprule
\textbf{Method} & \textbf{Path} & \textbf{Auth} & \textbf{Description} \\
\midrule
\endhead
\docpath{POST} & \docpath{/lessons} & Yes & Create lesson \\
\docpath{GET} & \docpath{/lessons} & Yes & List lessons \\
\docpath{GET} & \docpath{/lessons/grades} & Yes & List grades \\
\docpath{POST} & \docpath{/lessons/grades} & Yes & Create grade \\
\docpath{GET} & \docpath{/lessons/subjects} & Yes & List subjects \\
\docpath{POST} & \docpath{/lessons/subjects} & Yes & Create subject \\
\docpath{GET} & \docpath{/lessons/practice-available} & Yes & List lessons with templates \\
\docpath{POST} & \docpath{/lessons/:id/practice} & Yes & Start lesson practice \\
\docpath{GET} & \docpath{/lessons/:id} & Yes & Get lesson detail \\
\docpath{PUT} & \docpath{/lessons/:id} & Yes & Update lesson \\
\docpath{GET} & \docpath{/lessons/mastery/all} & Yes & Get current user mastery \\
\docpath{POST} & \docpath{/lessons/:id/study-time} & Yes & Track study time \\
\docpath{DELETE} & \docpath{/lessons/:id} & Yes & Delete lesson \\
\bottomrule
\end{longtable}

Practice start request:
\begin{lstlisting}[style=doccode]
{
  "difficulty": "easy"
}
\end{lstlisting}

Practice start response shape:
\begin{lstlisting}[style=doccode]
{
  "id": "attempt-uuid",
  "user_id": "user-uuid",
  "lesson_id": "lesson-uuid",
  "is_practice": true,
  "snapshots": [
    {
      "id": "snapshot-uuid",
      "generated_variables": {
        "a": 5,
        "b": 9
      },
      "right_answers": ["14"],
      "template": {
        "id": "template-uuid",
        "template_type": "arithmetic",
        "body_template_vi": "..."
      }
    }
  ]
}
\end{lstlisting}

\subsection{Tests API}

\begin{longtable}{L{0.12\textwidth}L{0.33\textwidth}L{0.12\textwidth}L{0.33\textwidth}}
\toprule
\textbf{Method} & \textbf{Path} & \textbf{Auth} & \textbf{Description} \\
\midrule
\endhead
\docpath{GET} & \docpath{/tests/grade-tests} & Yes & List grade tests \\
\docpath{POST} & \docpath{/tests/grade-tests/:gradeId/start} & Yes & Start grade test \\
\docpath{GET} & \docpath{/tests/attempts/:id} & Yes & Get attempt detail \\
\docpath{POST} & \docpath{/tests/submit-answer} & Yes & Submit answer for snapshot \\
\docpath{POST} & \docpath{/tests/attempts/:id/finish} & Yes & Finish attempt \\
\docpath{POST} & \docpath{/tests/templates} & Yes & Create one or many templates \\
\docpath{GET} & \docpath{/tests/templates} & Yes & List templates \\
\docpath{GET} & \docpath{/tests/templates/:id} & Yes & Get template detail \\
\docpath{PUT} & \docpath{/tests/templates/:id} & Yes & Update template \\
\docpath{DELETE} & \docpath{/tests/templates/:id} & Yes & Delete template \\
\docpath{POST} & \docpath{/tests/question-reports} & Yes & Report a question \\
\docpath{GET} & \docpath{/tests/question-reports} & Yes & List reports \\
\docpath{PATCH} & \docpath{/tests/question-reports/:id} & Yes & Update report status \\
\docpath{GET} & \docpath{/tests/my-practice-history} & Yes & List practice history \\
\bottomrule
\end{longtable}

Submit answer request:
\begin{lstlisting}[style=doccode]
{
  "snapshotId": "snapshot-uuid",
  "studentAnswer": "42"
}
\end{lstlisting}

Submit answer response:
\begin{lstlisting}[style=doccode]
{
  "isCorrect": true,
  "rightAnswers": ["42"],
  "explanation": "..."
}
\end{lstlisting}

Finish attempt response:
\begin{lstlisting}[style=doccode]
{
  "user_id": "user-uuid",
  "lesson_id": "lesson-uuid",
  "started_at": "2026-05-27T06:00:00.000Z",
  "total_score": 90,
  "newlyEarned": [
    {
      "slug": "first-step",
      "title_en": "First Step",
      "title_vi": "Buoc Dau Tien",
      "icon": "Footprints",
      "xp_reward": 50
    }
  ]
}
\end{lstlisting}

Question template payload:
\begin{lstlisting}[style=doccode]
{
  "lesson_id": "lesson-uuid",
  "template_type": "arithmetic",
  "difficulty": "medium",
  "body_template_en": "What is $a$ + $b$?",
  "body_template_vi": "$a$ + $b$ bang bao nhieu?",
  "explanation_template_en": "Add the two numbers.",
  "explanation_template_vi": "Cong hai so.",
  "logic_config": {
    "variables": {
      "a": { "min": 1, "max": 10 },
      "b": { "min": 1, "max": 10 }
    }
  },
  "accepted_formulas": ["a+b"]
}
\end{lstlisting}

\subsection{Achievements API}

\begin{longtable}{L{0.12\textwidth}L{0.25\textwidth}L{0.12\textwidth}L{0.41\textwidth}}
\toprule
\textbf{Method} & \textbf{Path} & \textbf{Auth} & \textbf{Description} \\
\midrule
\endhead
\docpath{POST} & \docpath{/achievements/seed} & Yes & Seed achievement definitions \\
\docpath{GET} & \docpath{/achievements} & Yes & List all achievements with earned state \\
\docpath{GET} & \docpath{/achievements/my} & Yes & List earned achievements \\
\docpath{POST} & \docpath{/achievements/check} & Yes & Trigger re-evaluation \\
\bottomrule
\end{longtable}

\subsection{Admin API}

\begin{longtable}{L{0.12\textwidth}L{0.30\textwidth}L{0.12\textwidth}L{0.36\textwidth}}
\toprule
\textbf{Method} & \textbf{Path} & \textbf{Auth} & \textbf{Description} \\
\midrule
\endhead
\docpath{GET} & \docpath{/admin/roles} & Yes & List basic roles \\
\docpath{GET} & \docpath{/admin/access-control} & Yes & List roles, actions, resources, subjects, and users \\
\docpath{POST} & \docpath{/admin/roles} & Yes & Create role \\
\docpath{PUT} & \docpath{/admin/roles/:id} & Yes & Update role \\
\docpath{PATCH} & \docpath{/admin/users/:id/role} & Yes & Assign role to user \\
\docpath{GET} & \docpath{/admin/users} & Yes & List users \\
\docpath{GET} & \docpath{/admin/users/:id/insights} & Yes & User insights \\
\docpath{POST} & \docpath{/admin/users} & Yes & Create user \\
\docpath{PATCH} & \docpath{/admin/users/:id} & Yes & Update user \\
\docpath{DELETE} & \docpath{/admin/users/:id} & Yes & Delete user \\
\docpath{GET} & \docpath{/admin/stats} & Yes & Dashboard statistics \\
\bottomrule
\end{longtable}

Role payload:
\begin{lstlisting}[style=doccode]
{
  "name": "teacher",
  "permissions": [
    { "action_id": 1, "resource_id": 2 },
    { "action_id": 2, "resource_id": 2 }
  ],
  "subject_ids": [1, 2]
}
\end{lstlisting}

\section{Sequence Flows}

\subsection{Login Flow}

\begin{enumerate}[leftmargin=1.5em]
  \item User submits email and password in the frontend.
  \item Frontend calls \docpath{POST /api/v1/auth/login}.
  \item Auth API finds the user by email together with role permissions.
  \item PostgreSQL returns the user record.
  \item Auth API validates the password with bcrypt.
  \item Auth API builds grouped permissions.
  \item Auth API signs the JWT.
  \item Frontend receives the token and user payload.
  \item Frontend saves the token in \docpath{localStorage} and a cookie.
\end{enumerate}

\subsection{Lesson Practice Flow}

\begin{enumerate}[leftmargin=1.5em]
  \item Student starts practice for a lesson.
  \item Frontend calls \docpath{POST /api/v1/lessons/:id/practice}.
  \item Lessons API validates the lesson and subject access.
  \item Lessons API loads question templates.
  \item For each selected template, the backend generates variables and evaluates accepted formulas through the Math Service.
  \item Backend creates a \docpath{test\_attempt}.
  \item Backend bulk-creates \docpath{question\_snapshots}.
  \item Backend returns the attempt and snapshot payload to the frontend.
\end{enumerate}

\subsection{Answer Submission Flow}

\begin{enumerate}[leftmargin=1.5em]
  \item Student submits an answer.
  \item Frontend calls \docpath{POST /api/v1/tests/submit-answer}.
  \item Tests API loads the \docpath{question\_snapshot} and template.
  \item For theoretical questions, the backend compares normalized text answers.
  \item For calculated questions, the backend calls the Math Service with \docpath{checkAnswer(formula, vars, answer)}.
  \item Backend updates the snapshot answer and result.
  \item Backend returns correctness, right answers, and explanation.
\end{enumerate}

\subsection{Finish Practice or Test Flow}

\begin{enumerate}[leftmargin=1.5em]
  \item Student finishes an attempt.
  \item Frontend calls \docpath{POST /api/v1/tests/attempts/:id/finish}.
  \item Tests API marks unanswered snapshots incorrect.
  \item Backend counts total and correct answers.
  \item Backend updates the \docpath{test\_attempt} score and completion state.
  \item For lesson-based practice, the backend calls the Mastery Service to upsert \docpath{user\_lesson\_mastery}.
  \item Backend calls the Level Service with \docpath{addXp(correct * 2, practice\_completion)}.
  \item Level Service inserts an \docpath{xp\_log} and updates \docpath{student\_stats}.
  \item Backend recomputes aggregate student stats.
  \item Backend calls the Achievement Service to evaluate and award achievements.
  \item Achievement Service may call the Level Service again to add XP for earned achievements.
  \item Backend returns score and newly earned achievements.
\end{enumerate}

\subsection{Admin Role Management Flow}

\begin{enumerate}[leftmargin=1.5em]
  \item Admin saves a role definition.
  \item Frontend calls \docpath{PUT /api/v1/admin/roles/:id}.
  \item Admin API updates the role name.
  \item Admin API deletes old role permissions.
  \item Admin API deletes old role subject permissions.
  \item Admin API inserts new role permissions.
  \item Admin API inserts new subject permissions.
  \item PostgreSQL returns the fresh role with relations.
  \item Backend returns the updated role payload.
\end{enumerate}

\section{Non-Functional Notes}

\subsection{Observed Design Strengths}

\begin{itemize}[leftmargin=1.5em]
  \item Clear frontend and backend separation
  \item Snapshot-based question persistence avoids re-generation drift
  \item Dynamic RBAC supports future role expansion
  \item Subject-level permissions provide scoped content access
  \item Progression pipeline is modularized into services
\end{itemize}

\subsection{Current Technical Constraints}

\begin{itemize}[leftmargin=1.5em]
  \item The frontend Prisma copy should mirror \docpath{backend/prisma/schema.prisma}; if both are kept, backend remains the canonical source of truth
  \item Template management permissions are enforced server-side with \docpath{manage:test}; frontend controls should remain aligned with that rule
  \item Achievement seeding is now opt-in through \docpath{AUTO\_SEED\_ACHIEVEMENTS=true} or the protected seed endpoint, so operational environments need an explicit seeding step
\end{itemize}

\section{Recommended Documentation Source of Truth}

For future maintenance, use these files as the implementation source of truth:
\begin{itemize}[leftmargin=1.5em]
  \item Backend API bootstrap: \docpath{backend/server.ts}
  \item Backend schema: \docpath{backend/prisma/schema.prisma}
  \item Authorization rules: \docpath{backend/middleware/auth.ts}
  \item Functional routes:
  \begin{itemize}[leftmargin=2em]
    \item \docpath{backend/routes/auth.ts}
    \item \docpath{backend/routes/lessons.ts}
    \item \docpath{backend/routes/tests.ts}
    \item \docpath{backend/routes/admin.ts}
    \item \docpath{backend/routes/achievements.ts}
  \end{itemize}
  \item Frontend API clients:
  \begin{itemize}[leftmargin=2em]
    \item \docpath{frontend/src/services/auth.ts}
    \item \docpath{frontend/src/services/lessonService.ts}
    \item \docpath{frontend/src/services/testService.ts}
    \item \docpath{frontend/src/services/adminService.ts}
    \item \docpath{frontend/src/services/achievementService.ts}
  \end{itemize}
\end{itemize}

\end{document}
